\section{K\"ahler differential structure}
\label{sec:kahler}

The shared-base tensor product has a corresponding decomposition of relative
differentials.

\begin{theorem}[Relative differential decomposition]
\label{thm:relative-differentials}
Let \(\cR\) be the iterated tensor product in
\cref{eq:quotient-tensor}.  There is a canonical \(\cR\)-module isomorphism
\begin{equation}
  \boxed{
  \Om_{\cR/A}\cong
  \bigoplus_{i=1}^{N}
  \left(
    \Om_{\cR^{(i)}/A}
    \otimes_{\cR^{(i)}}\cR
  \right).}
  \label{eq:relative-differential-decomposition}
\end{equation}
\end{theorem}

\begin{proof}
For two \(A\)-algebras \(B,C\), an \(A\)-derivation from \(B\otimes_A C\) to a
\((B\otimes_A C)\)-module \(M\) is uniquely determined by its restrictions to
\(B\) and \(C\).  Conversely, two such derivations define
\(d(b\otimes c)=c\,d_B(b)+b\,d_C(c)\).  The universal property of K\"ahler
differentials therefore gives
\begin{equation}
 \Om_{(B\otimes_A C)/A}
 \cong
 (\Om_{B/A}\otimes_B(B\otimes_A C))
 \oplus
 (\Om_{C/A}\otimes_C(B\otimes_A C)).
\end{equation}
Iterating this isomorphism and applying \cref{thm:quotient-tensor} proves the
claim.
\end{proof}

This theorem is global and requires no smoothness hypothesis.  Its meaning is
relative: after the shared parameter algebra is held fixed, different samples
have independent differential generators and relations.

\begin{proposition}[Absolute-to-relative transitivity]
\label{prop:transitivity}
For \(k\to A\to\cR\), there is a right-exact sequence
\begin{equation}
  \cR\otimes_A\Om_{A/k}
  \longrightarrow\Om_{\cR/k}
  \longrightarrow\Om_{\cR/A}
  \longrightarrow0.
  \label{eq:global-transitivity}
\end{equation}
\end{proposition}

\begin{proof}
This is the standard Jacobi--Zariski transitivity sequence for K\"ahler
differentials~\cite{StacksDifferentials}.  The first map sends a parameter
differential from \(A\) to its image in the global algebra; the second quotient
forgets parameter motion.
\end{proof}

The first term in \cref{eq:global-transitivity} is precisely the channel through
which sample-relative pieces couple in the absolute module.  Right exactness
does not assert that the first map is injective, nor that the sequence splits.

\subsection{Jacobian presentation}

Choose polynomial equality generators
\(F=(F_1,\ldots,F_m)\) for \(I\) in variables
\(x=(\theta,S^{(1)},\ldots,S^{(N)})\).  Then
\begin{equation}
  \cR^m\xrightarrow{J_F^{\mathsf T}}\cR^{P+\sum_iK_i}
  \longrightarrow\Om_{\cR/k}\longrightarrow0.
  \label{eq:global-jacobian-presentation}
\end{equation}
The image of \(J_F^{\mathsf T}\) records differential relations
\(dF_j=0\).  Different generating families can give different presentation
matrices while presenting isomorphic modules.

\begin{remark}[Global module versus fiber]
For a real solution \(p\), tensoring
\cref{eq:global-jacobian-presentation} with the residue field at \(p\) evaluates
all coefficients.  The resulting vector-space cokernel describes a cotangent
fiber.  Equality of such fibers, or equality of their row spaces, need not
imply equality of the original ideals or equality of their real projections.
\end{remark}

\begin{remark}[Singular points]
The module \(\Om_{\cR/k}\) and the exact sequence
\cref{eq:global-transitivity} exist at singular points.  What fails there is a
naive identification of module rank with the dimension of a smooth manifold.
Our algebraic decompositions remain valid; claims based on constant-rank
linear algebra require separate hypotheses.
\end{remark}
